\documentclass[11pt]{article}
\usepackage[margin=1in]{geometry}
\usepackage{amsmath,amssymb,amsthm,enumitem,hyperref}
\hypersetup{colorlinks=true,linkcolor=blue,urlcolor=blue,citecolor=blue}
\newtheorem{theorem}{Theorem}
\newcommand{\E}{\mathbb E}
\newcommand{\Q}{\mathcal Q}
\newcommand{\des}{\operatorname{des}}
\title{The word-enumerator formula for every arbor polytope\\\large MathDB 373577}
\author{Principia Math In Progress}
\date{September 18, 2026}
\begin{document}
\maketitle
\begin{abstract}
We prove Conjecture 5.2 of Athanasiadis, Xiao, and Yan for every arbor,
including arbitrary branching and nonsingleton vertex blocks. The result
is a direct corollary of the preorder descent theorem proved in our
companion MathDB 374958 solution. We give the precise arbor reduction
and reproduce the reciprocity and word-counting proof in full.
\end{abstract}

\section*{Task 1. Full solution}
An arbor $\tau$ on $[n]$ is a rooted tree whose vertices are the
nonempty blocks of a partition of $[n]$, where $n\ge1$.
Let $D(v)$ be the union of the blocks in the subtree rooted at $v$,
including its own block. The arbor polytope is
\[
 \Q_\tau=\{x\in\mathbb R_{\ge0}^n:x(D(v))\le |D(v)|
                             \text{ for every vertex $v$}\},
 \qquad x(A)=\sum_{i\in A}x_i.
\]
Let $W_\tau$ consist of words $w\in[n]^n$ in which the letters
from $D(v)$ occur at least $|D(v)|$ times for every $v$.
Let $\des(w)=\#\{j\in[n-1]:w_j>w_{j+1}\}$.
\begin{theorem}\label{thm:main}
For every arbor $\tau$ of size $n$,
\[
 \boxed{h^*(\Q_\tau,t)=\sum_{w\in W_\tau}t^{n-1-\des(w)}.}
\]
\end{theorem}

\subsection*{From descendant sets to order ideals}
Define a preorder $\sigma$ on $[n]$ by $i\le_\sigma j$ when
the block of $i$ is a descendant of the block of $j$. Elements
in one block are equivalent. Thus $D(v)$ is a principal order ideal.
Every order ideal $I$ of $\sigma$ is the disjoint union of the sets
$D(v)$ for its maximal blocks $v$. The subtrees are disjoint because
these blocks are incomparable in a rooted tree.

Summing over this disjoint union shows that the upper bounds on
$D(v)$ are equivalent to $x(I)\le |I|$ for every ideal $I$.
It follows that $\Q_\tau=\Q_\sigma$, where $\Q_\sigma$ denotes
the preorder polytope. The identical argument with lower bounds on
letter multiplicities gives $W_\tau=W_\sigma$, where $W_\sigma$
uses all order ideals. This reduction respects the usual order on
the alphabet, which need not be compatible with $\sigma$.

For completeness we prove the needed preorder formula below. Set $E=[n]$.
Write $\sigma^*$ for the reversed preorder and
$P_{\sigma^*}=\Q_{\sigma^*}\cap\mathbb Z^n$.

\subsection*{A counting identity from reciprocity}
We give the needed reciprocity argument in full. Let
\[
 \Q_\sigma(r,s)=\left\{x\in\mathbb R_{\ge0}^E:
       \sum_{e\in I}x_e\le r|I|+s
       \text{ for every nonempty order ideal }I\right\}.
\]
By \cite[Propositions 3.2--3.3 and 7.1]{AC}, for nonnegative
integers $r,s$ this is a lattice polytope, also expressible as an
integral Minkowski sum of simplices. Its lattice-point count is the
polynomial $\E_\sigma(r,s)$, of total degree $n$.

For positive integers $r,s$, the polytope is full dimensional
(it contains $[0,r]^E$). Its interior integer points satisfy
$x_e\ge1$ and $\sum_{e\in I}x_e\le r|I|+s-1$ for every
nonempty order ideal. Subtracting the all-ones vector therefore
bijects these points with $\Q_\sigma(r-1,s-1)\cap\mathbb Z^E$.
All inequality normals are nonzero; redundant inequalities have
no effect on this strict-interior characterization.

For $P=\Q_\sigma(r,s)$ one has $kP=\Q_\sigma(kr,ks)$ for
nonnegative integers $k$, so its Ehrhart polynomial equals
$\E_\sigma(kr,ks)$ as a polynomial in $k$. Ordinary Ehrhart
reciprocity at $k=-1$ and the translation just described imply
\begin{equation}\label{eq:rec}
 \E_\sigma(-u,-v)=(-1)^n\E_\sigma(u-1,v-1).
\end{equation}
Initially this holds for positive integer pairs $(u,v)$; it is a
polynomial identity because a polynomial vanishing on that entire
grid is zero (first fix $v$ and vary $u$, then vary $v$).

For $a=(a_e)_{e\in E}$ write $|a|=\sum_e a_e$ and $a_0=n-|a|$.
The established formula \cite[Lemma 4.1 and Equation (25)]{AC} is
\begin{equation}\label{eq:source}
 \E_\sigma(u,v)=\sum_{a\in P_{\sigma^*}}
       \binom{v+a_0}{a_0}\prod_{e\in E}\binom{u+a_e-1}{a_e}.
\end{equation}
Here $a_0\ge0$, and binomial coefficients with fixed nonnegative
lower argument are interpreted as polynomials in the upper argument.
Putting $(u,v)=(-m-1,-1)$ in \eqref{eq:source}, every term with
$a_0>0$ vanishes, since $\binom{a_0-1}{a_0}=0$.
For the remaining terms $|a|=n$, use
\[
 \binom{a_e-m-2}{a_e}=(-1)^{a_e}\binom{m+1}{a_e}.
\]
Together with \eqref{eq:rec}, this gives, for all $m\ge0$,
\begin{equation}\label{eq:count}
 \operatorname{Ehr}(\Q_\sigma,m)=\E_\sigma(m,0)
 =\sum_{\substack{a\in P_{\sigma^*}\\|a|=n}}
        \prod_{e\in E}\binom{m+1}{a_e}.
\end{equation}
This implication of bivariate reciprocity was already noted,
conditionally, in \cite[Remark 7.3]{AC}; equation \eqref{eq:rec}
supplies the missing premise.

\subsection*{An elementary multiset identity}
For a nonnegative integer vector $a$ with $|a|=n$, let $S_a$
be all words having $a_e$ copies of each $e$. We prove
\begin{equation}\label{eq:word}
 \sum_{m\ge0}\left(\prod_e\binom{m+1}{a_e}\right)t^m
 =\frac{\displaystyle\sum_{w\in S_a}t^{n-1-\des(w)}}{(1-t)^{n+1}}.
\end{equation}
Choose, for each letter $e$, a subset of $a_e$ bins from the
$m+1$ ordered bins $0,1,\ldots,m$. Place one copy of $e$ in each
chosen bin. The number of choices is the product on the left.
Within each bin, read its distinct letters in strictly decreasing
order, and concatenate the bins. This produces a word $w\in S_a$
and a weakly increasing sequence of bin indices
$0\le b_1\le\cdots\le b_n\le m$.

At each index $j$ for which $w_j\le w_{j+1}$, the bin index must
strictly increase. Conversely these conditions suffice: a constant
bin segment is strictly decreasing and contains no repeated letter.
Thus this construction is a bijection. The number of required
strict increases is $c(w)=n-1-\des(w)$. Subtract from $b_j$ the
number of required increases before position $j$. When $m\ge c(w)$
this gives a weakly increasing sequence of length $n$ in
$\{0,\ldots,m-c(w)\}$, counted by
$\binom{m-c(w)+n}{n}$. There are no such sequences for $m<c(w)$.
Consequently the contribution of $w$ to the generating function is
\[
 \sum_{m\ge c(w)}\binom{m-c(w)+n}{n}t^m
 =\frac{t^{c(w)}}{(1-t)^{n+1}},
\]
which proves \eqref{eq:word}, including repeated letters and $m=0$.

\begin{proof}[Proof of Theorem~\ref{thm:main}]
A vector $a$ of total sum $n$ belongs to $P_{\sigma^*}$ exactly
when $\sum_{e\in J}a_e\le |J|$ for every order filter $J$ of
$\sigma$. Taking complements $I=E\setminus J$, these conditions
are equivalent to $\sum_{e\in I}a_e\ge |I|$ for every order
ideal $I$ of $\sigma$. Therefore the sets $S_a$ indexed in
\eqref{eq:count} partition $W_\sigma$.
Sum \eqref{eq:word} over those vectors and use \eqref{eq:count}.
The definition
\[
 \sum_{m\ge0}\operatorname{Ehr}(\Q_\sigma,m)t^m
 =\frac{h^*(\Q_\sigma,t)}{(1-t)^{n+1}}
\]
now yields the desired identity.
\end{proof}

\subsection*{Verification and attribution}
The conjecture and the octopus special case are due to
Athanasiadis, Xiao, and Yan \cite{AXY}. The preorder counting formula
and the conditional deduction from bivariate reciprocity are due to
Athanasiadis and Chapoton \cite{AC}. The complete reciprocity proof
above is reproduced from our MathDB 374965 and 374958 packages;
this separately listed arbor target is a corollary, not an independent
new general theorem. Since $\Q_\tau=\Q_\sigma$ and
$W_\tau=W_\sigma$, the formula proved above is exactly the target.

All 14 pages of the current arbor primary and all 26 pages of the
preorder paper were read. The verifier checks every one of the
1,674 labeled arbors through size five, including arbitrary blocks:
11,275 ideal decompositions, 1,674 complete $h^*$-identities,
9,866 binomial counts, and 740 direct lattice enumerations.
It also checks the published octopus formula and six-element examples.
Bounded checks found no earlier external full target claim or identified
target journal final. Review is by the authoring agent only.

\section*{Tasks 2--5. Proposed follow-up directions}
\begin{enumerate}[label=\textbf{Task \arabic*.},start=2,leftmargin=*]
\item Construct a recursive bijection from dilated arbor lattice
points to the bin assignments, using the branching of the rooted tree.
\item Refine the formula by descent sets and derive shuffle identities
for the subtrees attached to a vertex.
\item Find a parking interpretation of the magic-basis coefficients
for arbitrary arbors, extending the known octopus case.
\item Determine which modifications of the subtree capacities admit
an analogous word formula with a fixed multiplicity constraint.
\end{enumerate}
\begin{thebibliography}{9}
\raggedright
\bibitem{AXY} C. A. Athanasiadis, Q. Xiao, and X. Yan,
\emph{Lattice point enumeration of some arbor polytopes},
\href{https://arxiv.org/abs/2603.11654v1}{arXiv:2603.11654v1}, 2026.
\bibitem{AC} C. A. Athanasiadis and F. Chapoton,
\emph{Polytopes and posets associated to preorders},
\href{https://arxiv.org/abs/2605.26916v1}{arXiv:2605.26916v1}, 2026.
\end{thebibliography}
\end{document}
